\section{Resultados de la implementación}

\subsection{Centralización y seguridad de la información del pasajero}

Uno de los resultados más críticos ha sido la transición de la gestión de datos sensibles desde canales informales como WhatsApp y Excel hacia la ficha centralizada del cliente en Odoo.

\begin{itemize}
    \item[a.] \textbf{Gestión segura de documentos}: se ha eliminado la práctica de recibir fotografías de pasaportes por medios inseguros y borrarlas periódicamente por temor a vulnerabilidades. Ahora, los documentos se almacenan en el repositorio estándar de Odoo, garantizando la integridad de la información para futuras consultas o emergencias durante el tour.
    \item[b.] \textbf{Visión 360° del turista}: la gerencia dispone ahora de un historial completo de interacciones, preferencias y requisitos específicos del cliente extranjero, lo que facilita una atención personalizada y profesional.
\end{itemize}

\subsection{Optimización del proceso comercial y seguimiento}

El ciclo de venta ha sido profesionalizado mediante la implementación del embudo de ventas (pipeline).

\begin{itemize}
    \item[a.] \textbf{Trazabilidad de oportunidades}: se ha logrado un seguimiento sistemático de los prospectos provenientes de recomendaciones y eventos, evitando que las consultas se pierdan en hilos de chat individuales.
    \item[b.] \textbf{Estandarización de cotizaciones}: el uso de plantillas nativas permite generar propuestas comerciales coherentes con la imagen de lujo de la empresa, integrando automáticamente términos y condiciones legales sobre cancelaciones y penalidades.
\end{itemize}

\subsection{Control operativo de itinerarios complejos}

La configuración del módulo de Proyectos bajo una vista Kanban ha permitido ordenar la logística de viajes que varían entre 2 y 20 días.

\begin{itemize}
    \item[a.] \textbf{Hitos de servicio}: se visualiza claramente la etapa en la que se encuentra cada turista (desde el borrador de cotización hasta la ejecución final), facilitando la coordinación oportuna con proveedores de hoteles y transporte.
    \item[b.] \textbf{Gestión documental operativa}: los boletos de tren, avión y entradas a atractivos se cargan directamente en la tarea correspondiente, reduciendo errores de digitación que anteriormente ponían en riesgo el ingreso de los pasajeros debido a las estrictas regulaciones de los sitios arqueológicos.
\end{itemize}

\subsection{Mejora en el control financiero de reservas}

Se ha establecido un flujo riguroso para la gestión de pagos parciales, que usualmente oscilan entre el 30\% y 50\%.

\begin{itemize}
    \item[a.] \textbf{Visibilidad de saldos}: el sistema permite identificar de forma inmediata si un cliente tiene saldos pendientes antes de proceder con compras no reembolsables de boletos o servicios, mitigando riesgos financieros para la agencia.
    \item[b.] \textbf{Registro automatizado}: el registro de abonos dentro de Odoo genera un comprobante interno profesional que refuerza la confianza del turista extranjero en la legalidad de la empresa.
\end{itemize}

\subsection{Indicadores de desempeño (KPIs)}

Para validar el éxito de la implementación, se han proyectado los siguientes indicadores basados en la mejora operativa:

\begin{table}[h]
    \centering
    \begin{tabularx}{\textwidth}{|l|l|l|X|}
        \hline
        \textbf{Indicador (KPI)} & \textbf{Línea base (manual)} & \textbf{Meta esperada (Odoo)} & \textbf{Impacto} \\
        \hline
        Tiempo de respuesta al cliente & 2--4 horas (promedio) & $<$ 1 hora & -50\% a -70\% de reducción \\
        \hline
        Conversión de cotizaciones & Sin registro sistemático & +15\% a +20\% & Incremento en ventas \\
        \hline
        Seguridad de pasaportes & 0\% (almacenamiento volátil) & 100\% (cifrado en la nube) & Mitigación de riesgos \\
        \hline
        Visibilidad de saldos & Consulta manual en Excel & Acceso en tiempo real & Control interno crítico \\
        \hline
    \end{tabularx}
    \caption{Indicadores clave de desempeño tras la implementación}
\end{table}

\subsection{Impacto en la eficiencia operativa}

Globalmente, la implementación ha reducido la carga administrativa de la gerencia en aproximadamente un 20\%, permitiendo delegar tareas operativas y enfocarse en la expansión hacia los mercados internacionales. La eliminación de la doble digitación y la centralización de la data han sentado las bases para un crecimiento proyectado del 30\% en las reservas digitales directas.

